\inicializarSubseccion{Analisis de la séptima pregunta (Por parte de Andrés)}
\textbf{7. De los siguientes rangos de precios, ¿Cuál consideras más adecuado para el adaptador?}

\tab Esta pregunta tiene estructura de conjunto; dando al usuario un rango de precios en donde el valor del producto está definido. Nuestra tabla de elementos quedaría algo tal que así

\begin{equation}
    A = [100, 150], B = [150, 200], C = [200, 250]
\end{equation}

\tab Con estos datos, es fácil ver que nuestro universo está definido como la cantidad de usuarios que contestaron la encuesta ($|U|=102$) y que la cantidad de elementos que conforman el universo $|A|=22, |B|=62, |C|=18$. Utilizando operadores de diferencia, es fácil ver que son completamente independientes. Es decir que:

\begin{equation}
    A \not\subset B \not\subset C \to (A\cap B),(A\cap C),(B\cap C) = \emptyset
\end{equation}

\tab Gracias a esto, sabemos que la moda de precios es absoluta, por ende sabemos que el rango B al tener la mayor cantidad de elementos (el 61.4\% de los datos). Lo que quiere decir que la gran mayoria de personas estan dispuestas a pagar por un producto de calidad pero a un precio accesible.